
%%%%%%%%%%%%%%%%%%%%%%%%%%%%%%%%%%%%%%%%%%%%%%%%%%%%%%%%%%%%%
\subsubsection{具体分析}

问题四要求将排放标准从$10\;\text{mg/Nm}^3$收紧至$5\;\text{mg/Nm}^3$，定量分析各工况电耗增幅并给出高浓度工况的应对建议，属于约束收紧影响评估类问题。

\textbf{重求解vs灵敏度外推。}最直接的想法是用问题三的灵敏度外推——已知$C_{limit}$从10降到5，用$\Delta P\approx S^P\cdot\Delta x$估算。但这是线性近似，$C_{limit}$减半是大幅变动，线性外推误差大。更准确的方案是直接将$C_{limit}'=5$代入问题二的优化模型重新求解。因代码中\texttt{solve\_one\_regime(...,\,C\_limit)}已将排放阈值参数化，问题四只需传入$C_{limit}'=5$即可复用全部寻优逻辑。

\textbf{电耗增幅量化。}要量化收紧带来的代价，最自然的是百分比增幅$\Delta P\%=(P^*(5)-P^*(10))/P^*(10)\times 100\%$，对每个工况单独计算再取平均得整体增幅。

\textbf{可行性校验。}收紧到$5\;\text{mg/Nm}^3$后，某些工况可能即使电压推到$U_{max,i}$也无法满足$C_{out}\leq 5$，此时解在物理可行域外，需标记并提示"需硬件改造（增加电场数或提升变压器容量）"。
